\documentclass[a4paper]{resume} % Use the custom resume.cls style
\usepackage{nopageno}
\usepackage[left=0.25in,top=0.15in,right=0.25in,bottom=0.3in]{geometry} % Document margins
\usepackage{fontawesome}
\usepackage[utf8]{inputenc}
\usepackage{times}
    \newcommand{\tab}[1]{\hspace{.2667\textwidth}\rlap{#1}}
    \newcommand{\itab}[1]{\hspace{0em}\rlap{#1}}
%\usepackage[german]{babel}
\usepackage{ulem}
\usepackage{enumitem}
    \setitemize{noitemsep,topsep=0pt,parsep=0pt,partopsep=0pt}
\usepackage {graphicx}

\name{José Alfredo Rosas Córdova} % Your name 
\address{\faEnvelope{ jose.alfredo.rosas.cordova@stud.uni-hannover.de}\hspace{0.75em}\faPhone{ +49 162 519 7979}}
\address{\faMapMarker{ Heidjerhof 3, Zi. 103, 30625 Hannover}} % Your address
\address{\faUser{ October 22, 1999}\hspace{0.75em}\faGlobe{ Guadalajara, Mexico}}

\begin{document}
    % === AUSBILDUNG ===
    \begin{rSection}{Education}
    	{\bf Mechanical Engineering, M. Sc.} \hfill {\itshape 10/2025 --- today} 
    	\\{ \textit{Leibniz University Hannover}}
    	%\\ Notendurchschnitt: 97.97 (entspricht 1.2)    	
    	
        {\bf Bachelor of Aeronautical Engineering} \hfill {\itshape 08/2017 --- 06/2023} 
        \\{ \textit{Universidad Autónoma de Nuevo León}}
        \\ Graduated summa cum laude
        
        {\bf Combined Study and Practice Stays (DAAD-KOSPIE)} \hfill {\itshape 09/2021 --- 09/2022} 
        \\{\bf Mechanical Engineering, B. Sc.}
        \\ { \textit {Technische Universität Braunschweig}}
        \\ GPA: 1.9

 
        
    \end{rSection}
    
    \begin{rSection}{Languages}
        \begin{tabular}{ @{} >{\bfseries}l @{\hspace{6ex}} l }
            Spanish & Native speaker\\
            English & TOEFL ITP -- 677 Points (C1)\\
            German & TestDaF -- TDN 5 (C1)
        \end{tabular}
    \end{rSection}
    
    \begin{rSection}{IT Skills}
        \begin{tabular}{ @{} >{\bfseries}l @{\hspace{6ex}} l }
            CAD \& CAM: & PTC Creo, OnShape, SOLIDWORKS, Fusion 360 \\
            Programming: & Python, MATLAB \\
            Simulation \& Modelling: & ANSYS Workbench, Fluent, Simulink, Xfoil\\
            Others: & MS Office, Oracle PLM \& ERP, \LaTeX, Git\\
        \end{tabular}
    \end{rSection}

    \begin{rSection}{Work Experience}
   		{\bf Research Assistant} \hfill{\itshape 04/2026 --- today} 
   		\\{\textit{\textbf{Institut für Turbomaschinen und Fluid-Dynamik} -- Hannover, Germany}}
   		\begin{itemize}
   			\item Use of specialized software for the extraction of deformation information of a wind turbine with marked blades through the use of digital image correlation techniques.
   			\item Post-processing of obtained large-volume datasets using self- and further-developed scripts using Python and MATLAB relying heavily on parallel programming schemes. The processing steps included but were not limited to classification, transformation and mapping of the measured points as well as the generation of diagrams needed for the subsequent validation steps.
   			\item Heavy use of Git-based tools to manage software versions and the safe implementation of new functionalities without putting the current production version at risk.
   			\item Use of a Blender-based digital twin to validate the correctness of the implemented algorithms before using them to evaluate real-world experimental data.
   			
   		\end{itemize}
   		
    	
    	
        {\bf Mechanical Engineer - New Product Development \& Introduction, Liquid Cooling} \hfill{\itshape 06/2023 --- 09/2025} 
        \\{\textit{\textbf{Vertiv} -- Monterrey, Mexico}}
        \begin{itemize}
            \item Development, documentation and sustaining of mechanical components and assemblies encompassing sheet metal welded structures; primary and secondary cooling circuits (i.e. refrigerant and water, copper and stainless steel based piping networks, respectively) for new products tailored towards high density data center cooling applications (e.g. chillers, CDUs, secondary fluid networks). 
            \item Experience with design development within a collaborative work environment based on the use of PLM and ERP systems based on the use of Engineering Change Orders and strict revision control.
            \item Execution of both DFM\&A as well as DFMEA studies during the design and development process of new products for the identification and correction of areas of opportunity before the parts are committed to production.
            \item Participation in the design validation testing for both complete, production-ready units as well as engineering prototypes for both performance assessments and certification by external entities (e.g. UL, NOM).
            \item Collaboration with manufacturing, field teams and facilities for the implementation of design improvements around manufacturing as well as serviceability concerns within both development and production environments.
            \item Participation in the planning, part procurement, and build process of both engineering, and manufacturing prototypes of brand new products both locally and abroad.
        \end{itemize}
        
        {\bf Project Engineer - UAV Development} \hfill{\itshape 01/2023 --- 04/2023} 
        \\{\textit{\textbf{Spacelab MX} -- Monterrey, Mexico}}
        \begin{itemize}
            \item Initial aircraft sizing and preliminary flight performance studies.
            \item Design of the initial airframe structural and drivetrain components for the prototype phase of the project.
            \item Integration, ground-, and initial flight testing of the testbed aircraft and its software.
        \end{itemize}
    
        {\bf System Testing \& Mechanical Design Intern - Engineering Mechatronics} \hfill{\itshape 04/2022 --- 08/2022} 
        \\{\textit{\textbf{Bosch Rexroth AG} -- Stuttgart, Germany}}
        \begin{itemize}
            \item Design and implementation of system tests within a production environment for autonomous mobile robots (UGV) based around previously defined KPIs established around beta-tester customer requirements.
            \item Evaluation, analysis, and fault detection of the drive performance of the UGVs as well as the communication layer between UGVs and fleet management service based on test data.
            \item Handling and maintenance of autonomous mobile robots.
            \item Mechanical design of internal hardware elements and drive mechanisms.
        \end{itemize}
    \end{rSection}

    
    \begin{rSection}{Research Experience}
        {\bf Research Assistant - Aerodynamics Laboratory} \hfill {\itshape 07/2020 --- 06/2023} 
        \\{\textit{\textbf{Centro de Investigación e Innovación en Ingeniería Aeronáutica} -- Monterrey, Mexico}}
        \vspace{-0.5em}
        \begin{itemize}
            \item Co-published design case study paper on \textit{Aircraft Engineering and Aerospace Technology} journal.\hfill {\itshape 09/2025}
            \begin{itemize}
                \item \textit{Design, manufacture and testing of a low-cost tilt-rotor UAV: a practical development case} (DOI: 10.1108/AEAT-05-2025-0168). 
            \end{itemize}
            
            \item Co-published potential flow solver validation paper on the \textit{Revista Internacional de Métodos Numéricos para Cálculo y Diseño en Ingeniería} journal.\hfill {\itshape 11/2024}
            \begin{itemize}
            	\item \textit{Validation of VSPAERO for Basic Wing Simulation} (DOI: 10.23967/j.rimni.2024.10.56492). 
            \end{itemize}            
        \end{itemize}
    \end{rSection}
    
    
    
    % === Extracurricular Section ===
    \begin{rSection}{Extracurricular Activities}
         {\bf Club Member} \hfill {\itshape 10/2025 --- today} 
         \\{\textit{\textbf{Akademischer Fliegergruppe Hannover e.V.}}}
         \vspace{-0.5em}
         \begin{itemize}
         	\item Participation in the maintenance of the club facilities and aircraft.
         \end{itemize}
        
        {\bf Team Member (Technical Report Lead, Chief Engineer)} \hfill {\itshape 09/2019 --- 06/2022} 
        \\{\textit{\textbf{Axios Aero Design}}}
        \vspace{-0.5em}
        \begin{itemize}
        	\item Technical report production consisting primarily of a documentation of the engineering design process as well as technical drawings describing the final product to be presented in the yearly SAE competition.
            \item Aircraft conceptual and preliminary design studies encompassing aerodynamics, structures, electric propulsion, and flight stability and control.
            \item Project planning and execution based around system engineering methodologies as well as assistance in the detailed design of structural aircraft components such as wing, fuselage, and empennage structures.
        \end{itemize}
    \end{rSection}
    
    % === PRICES ===
    \begin{rSection}{Awards}
        {\bf SAE Aero Design México 2022} \hfill {\itshape 03/2022}
        \\{\textit{2. Place Overall}}
        \\{\textit{3. Place Technical Report}}
    
        {\bf MathWorks Minidrone Competition Mexico 2021} \hfill {\itshape 08/2021}
        \\{\textit{1. Place}}
        
        % {\bf SAE Aero Design México 2020} \hfill {\itshape 10/2020}
        % \\{\textit{6. Place overall}}
        % \\{\textit{4. Place Technical Report}}
    \end{rSection}

    % === SIGNATURE ===
    \begin{rSection}{}
        \vspace{7.5em}
        \begin{minipage}{0.5\linewidth}
            \uline{San Nicolás de los Garza, NL \today}\\
            \textit{Place, Date}
        \end{minipage}
        \begin{minipage}{0.5\linewidth}
            \uline{%
            	José Alfredo Rosas Córdova%
            	\hspace{0.em}\rlap{\raisebox{-1.75em}[0pt][0pt]{\includegraphics[height=7em]{/home/jarc/Nextcloud/100 Personal/110 Documents/Recursos Gráficos/blue_signature.png}}}%
            	\hspace{5em}%
            }\\
            \textit{Signature}
        \end{minipage}
    \end{rSection}
\end{document}